% Part: normal-modal-logic
% Chapter: syntax-and-semantics
% Section: introduction

\documentclass[../../../include/open-logic-section]{subfiles}

\begin{document}

\olfileid[hi]{nml}{syn}{int}

\olsection{परिचय}

मोडल तर्कशास्त्र \emph{मोडल प्रतिज्ञप्तियों} और उनके बीच के परिणाम संबंधों
का अध्ययन करता है। मोडल प्रतिज्ञप्तियों के उदाहरण ये हैं:
\begin{enumerate}
\item यह आवश्यक है कि $2+2=4$।
\item यह आवश्यक रूप से संभव है कि कल वर्षा होगी।
\item यदि यह आवश्यक रूप से संभव है कि~$!A$, तो यह संभव है कि~$!A$।
\end{enumerate}
संभावना और आवश्यकता ही एकमात्र मोडलताएँ नहीं हैं: अन्य एकपदी संयोजकों को भी
मोडलता माना जाता है, जैसे ``ऐसा होना चाहिए कि~$!A$'', ``ऐसा होगा
कि~$!A$'', ``डाना जानती है कि~$!A$'', या ``डाना मानती है कि~$!A$''।

मोडल तर्कशास्त्र सबसे पहले अरस्तू की \emph{De Interpretatione} में दिखाई
देता है। उन्होंने पहली बार ध्यान दिया कि आवश्यकता से संभावना निकलती है,
पर इसका विलोम नहीं; संभावना और आवश्यकता को एक-दूसरे से परिभाषित किया जा
सकता है; यदि $!A \land !B$ का सत्य होना संभव है, तो $!A$ का सत्य होना
संभव है और $!B$ का सत्य होना संभव है, पर इसका विलोम नहीं; और यदि
$!A \to !B$ आवश्यक है, तो $!A$ आवश्यक होने पर $!B$ भी आवश्यक है।

मोडल तर्कशास्त्र का पहला आधुनिक दृष्टिकोण सी.~आई. लुईस का कार्य था, जिसकी
परिणति लुईस और लैंगफ़र्ड की \emph{Symbolic Logic} (1932) में हुई। लुईस और
लैंगफ़र्ड भौतिक सशर्त द्वारा निहितार्थ के निरूपण से संतुष्ट नहीं थे:
$!A \lif !B$, ``$!A$ से $!B$ निकलता है'' का खराब विकल्प है। इसके स्थान पर
उन्होंने निहितार्थ को ``आवश्यक रूप से, यदि $!A$ तो $!B$'' के रूप में
अभिलक्षित करने का प्रस्ताव रखा, जिसे $!A \strictif !B$ से चिह्नित किया।
भिन्न गुणों को व्यवस्थित करने के प्रयास में लुईस ने पाँच भिन्न मोडल तंत्र,
$\Log{S1}, \ldots, \Log{S4}, \Log{S5}$, पहचाने; इनमें अंतिम दो आज भी
प्रयुक्त होते हैं।

लुईस और लैंगफ़र्ड का दृष्टिकोण पूर्णतः \emph{वाक्यविन्यासगत} था: उन्होंने
युक्तिसंगत अभिगृहीत और नियम पहचाने और जाँचा कि उन साधनों से क्या सिद्ध किया
जा सकता है। अर्थगत दृष्टिकोण बहुत समय तक हाथ नहीं आया। अंततः रुडोल्फ़
कार्नाप ने \emph{Meaning and Necessity} (1947) में \emph{अवस्था-वर्णन} की
धारणा का प्रयोग करके पहला प्रयास किया। अवस्था-वर्णन परमाण्विक वाक्यों का
एक संग्रह है (वे वाक्य जो उस अवस्था-वर्णन में ``सत्य'' हैं)। सत्य की
परिभाषा को स्वेच्छ वाक्यों $!A$ तक विस्तृत करने के बाद कार्नाप $!A$ को
\emph{आवश्यक रूप से सत्य} तब परिभाषित करते हैं जब वह सभी अवस्था-वर्णनों में
सत्य हो। कार्नाप का दृष्टिकोण \emph{पुनरावृत्त} मोडलताओं को नहीं सँभाल
सकता था, क्योंकि ``संभवतः आवश्यक रूप से \ldots संभवतः $!A$'' रूप के वाक्य
सदैव सबसे भीतरी मोडलता तक घट जाते हैं।

मोडल अर्थविज्ञान में बड़ी सफलता सॉल क्रिप्के के लेख ``A Completeness Theorem
in Modal Logic'' (JSL 1959) से मिली। क्रिप्के ने अपना कार्य लाइब्नित्स के इस
विचार पर आधारित किया कि कोई कथन आवश्यक रूप से सत्य है यदि वह ``सभी संभव
जगतों में'' सत्य हो। किंतु यह विचार कार्नाप के विचार जैसी ही कमी से ग्रस्त
है, क्योंकि किसी जगत $w$ (या अवस्था-वर्णन $s$) में कथन की सत्यता $w$ पर
बिल्कुल निर्भर नहीं होती। इसलिए क्रिप्के ने माना कि जगत एक
\emph{अभिगम्यता संबंध} $R$ से जुड़े होते हैं और ``आवश्यक रूप से $!A$'' रूप
का कथन किसी जगत $w$ में तभी सत्य है जब $!A$, $w$ से \emph{अभिगम्य} सभी
जगतों $w'$ में सत्य हो। इस दृष्टिकोण का कोई रूप देने वाले अर्थविज्ञान को
क्रिप्के अर्थविज्ञान कहते हैं; उसने मोडल तर्कशास्त्रों (बहुवचन में) के तीव्र
विकास को संभव बनाया।

क्रिप्के अर्थविज्ञान से व्याख्यायित होने पर मोडल तर्कशास्त्र हमें दिखाता है
कि \emph{संबंधात्मक संरचनाएँ} ``भीतर से'' कैसी दिखती हैं। संबंधात्मक संरचना
केवल एक द्विपदी संबंध से युक्त समुच्चय है (उदाहरणार्थ, सामाजिक सुरक्षा
संख्या के क्रम में रखे गए कक्षा के विद्यार्थियों का समुच्चय एक संबंधात्मक
संरचना है)। वास्तव में संबंधात्मक संरचनाओं के प्रक्षेत्र अनेक प्रकार के
होते हैं: जगत की अवस्थाओं की सापेक्ष संभावना के अतिरिक्त किसी अभिकर्ता की
ज्ञानमीमांसात्मक संभावना से जुड़ी ज्ञानमीमांसात्मक अवस्थाएँ, या अवस्था-
संक्रमणों से जुड़ी किसी गतिक तंत्र की अवस्थाएँ आदि हो सकती हैं। इन सभी का
प्रतिरूपण मोडल तर्कशास्त्र से किया जा सकता है: पहला हमें साधारण एलेथिक मोडल
तर्कशास्त्र देता है; अन्य हमें ज्ञानमीमांसात्मक तर्कशास्त्र, गतिक
तर्कशास्त्र आदि देते हैं।

हम एक विशेष दृष्टिकोण पर ध्यान देंगे, जिसे मोडल तर्कशास्त्री ``अनुरूपता
सिद्धांत'' कहते हैं। क्रिप्के की सबसे महत्त्वपूर्ण आरंभिक खोजों में से एक यह
थी कि अभिगम्यता संबंध~$R$ के अनेक गुणों (जैसे उसका संक्रमणशील या सममित होना)
को उपयुक्त ``मोडल स्कीमाओं'' द्वारा स्वयं \emph{मोडल भाषा में} अभिलक्षित
किया जा सकता है। उदाहरणार्थ, मोडल तर्कशास्त्री कहते हैं कि $R$ की
स्वपरावर्तकता इस स्कीमा के ``अनुरूप'' है: ``यदि आवश्यक रूप से~$!A$, तो
$!A$''। हम मुख्यतः मोडल तर्कशास्त्र के उन शास्त्रीय तंत्रों (जैसे $\Log{S4}$
और $\Log{S5}$) के अनुरूपता सिद्धांत का अध्ययन करेंगे जो स्कीमाओं $\Ax{D}$,
$\Ax{T}$, $\Ax{B}$, $\Ax{4}$ और~$\Ax{5}$ के संयोजन से प्राप्त होते हैं।

\end{document}
